\documentclass[12pt]{article}
\usepackage[margin=1in]{geometry}
\usepackage{multirow}
\usepackage{forest}

\begin{document}

\noindent
\textbf{Name:}\underline{\hspace{6cm}} \hfill \textbf{Date:}\underline{\hspace{3cm}}

\vspace{0.5em}
\noindent\textit{Side 1 of 2}

\vspace{1.5em}
\noindent\textbf{1. Consider the extensive-form game below. Player 1 moves first, choosing $L$ or $R$. If Player 1 chooses $R$, Player 2 then chooses $U$ or $D$. Payoffs are listed as (Player 1's payoff, Player 2's payoff).}

\begin{center}
\begin{forest}
[$\varepsilon$
    [{$(L,*)$}
        [{$(4,2)$}]
    ]
    [{$(R,*)$}
        [{$(*,U)$}
            [{$(1,1)$}]
        ]
        [{$(*,D)$}
            [{$(3,5)$}]
        ]
    ]
]
\end{forest}
\end{center}

\noindent \textbf{State the full subgame perfect equilibrium (as an action plan for each player) and the resulting equilibrium payoffs.}

\vfill
\newpage

\noindent\textit{Side 2 of 2}

\vspace{1.5em}
\noindent \textbf{2. Please summarize the SPE existence proof as asked on the homework. Your summary should clearly explain \#1) How the base case is set up for the induction, \#2) Where the inductive hypothesis is used, \#3) How the candidate SPE is constructed, \#4) Where and why Nash's existence result is used (it is used multiple times), \#5) Where and why the one shot deviation principle is used, and \#6) The ultimate source of the final contradiction.}

\vfill

\end{document}
